\section{Hormone Routing}
\label{sec:hormone}

Hormone routing injects a learned linear combination of frozen behavioural direction vectors into the residual stream after every attention block. This section presents the routing module, the contrast-pair extraction protocol used to obtain the directions, and the staged training schedule by which the directions are extracted, frozen, and re-injected into pretraining.

\subsection{The hormone routing module}

Let $V \in \mathbb{R}^{\nhormones \times \dmodel}$ be a fixed bank of $\nhormones$ unit-norm direction vectors, each of dimension $\dmodel$, attached as a non-trainable buffer to the model. We refer to these as \emph{hormone vectors}; their derivation is described in Section~\ref{subsec:extraction}. The hormone routing module operates per residual position, treating each token's hidden state $h \in \mathbb{R}^{\dmodel}$ independently.

\paragraph{Per-token router.} Given a hidden state $h$ at the position where the module is applied, we compute a token-conditional convex combination over the hormone bank:
\begin{equation}
    w(h) = \mathrm{softmax}(W_R \, h) \in \Delta^{\nhormones - 1}, \qquad W_R \in \mathbb{R}^{\nhormones \times \dmodel}.
\end{equation}
The router $W_R$ is a single linear map without bias and the only routing parameter introduced. Because $\nhormones$ is small (we use $\nhormones = 7$ throughout), the router adds only $\nhormones \dmodel$ parameters per attention layer.

\paragraph{Per-hormone magnitudes.} Each hormone has a single learned scalar magnitude $m_i$, collected as a vector $m \in \mathbb{R}^{\nhormones}$ initialised to all-ones. The magnitudes scale the corresponding rows of the hormone bank:
\begin{equation}
    \tilde V = m \odot V.
\end{equation}
This decouples the geometric direction of each hormone (kept frozen and unit-norm) from its overall expressive strength (free to grow or shrink during training). Magnitudes are excluded from weight decay.

\paragraph{Output gate.} The router output is computed as
\begin{equation}
    \Delta h = g \cdot \big( w(h)^\top \, \tilde V \big),
\end{equation}
where $g \in \mathbb{R}$ is a single scalar gate, initialised to zero, multiplied uniformly across all tokens and all hormones. At initialisation $\Delta h \equiv 0$ and the model is exactly the vanilla hybrid backbone; the gate is excluded from weight decay so that early-training gradients can drive it freely. The hormone shift is added to the residual stream after the MLP residual addition (Section~\ref{sec:arch}).

\paragraph{Total parameter cost.} Per attention layer: $\nhormones \dmodel + \nhormones + 1$ parameters (router, magnitudes, gate). The hormone bank $V$ itself is shared across all attention layers, contributing $\nhormones \dmodel$ non-trainable parameters once. For \textsc{small} ($\dmodel = 1536$, $\nhormones = 7$, three attention layers): $\approx 32{,}000$ trainable parameters total, under $0.01\%$ of model parameters.

\subsection{Document-style contrast-pair extraction}
\label{subsec:extraction}

The hormone vectors are obtained by passing pairs of contrasting documents through the partially-trained base model and computing the difference of mean residual-stream activations across the two pair members.

\paragraph{Why document style.} Existing activation-steering work~\citep{turner2023activation,rimsky2024steering} typically uses \emph{instruction-style} contrast pairs --- e.g.\ a system prompt instructing the model to be honest versus to be deceptive --- because the target models are instruction-tuned. In the pretraining setting these prompts are out-of-distribution: the base model has never seen ``You are a helpful assistant'' as a prefix and its activations on such inputs reflect the prompt's surface form more than any underlying behavioural axis.

We instead use pairs of \emph{naturally occurring document genres} that span the same affective or stylistic dimension while differing in content. For instance, an emergency-dispatch transcript and a leisurely fishing-trip narrative both describe ongoing events, but they differ sharply in urgency. By pairing seven such contrasts (urgency, formality, certainty, warmth, curiosity, caution, decisiveness), each instantiated as multiple pairs across deliberately varied genres, we recover difference-of-means directions that capture the affective dimension itself rather than any single genre's surface vocabulary.

\paragraph{Extraction procedure.} For each contrast (pair of document sets) $(\mathcal{D}^+, \mathcal{D}^-)$ associated with a hormone $i$, and for each pair $(d^+, d^-) \in \mathcal{D}^+ \times \mathcal{D}^-$, we run forward passes of the base model and collect the residual-stream activations after every attention block, averaging over both attention layers and token positions to obtain $\bar h^+(d^+), \bar h^-(d^-) \in \mathbb{R}^{\dmodel}$. The hormone direction is the L2-normalised difference of means across pairs:
\begin{equation}
    v_i = \mathrm{normalize}\!\left( \frac{1}{|\mathcal{D}^+|} \sum_{d^+} \bar h^+(d^+) - \frac{1}{|\mathcal{D}^-|} \sum_{d^-} \bar h^-(d^-) \right).
\end{equation}
Stacking $\{v_i\}_{i=1}^{\nhormones}$ produces the hormone bank $V$.

\paragraph{Why average over attention layers.} The residual stream is shared across all layers but the contrast signal is sharpest at the layers where the relevant decisions are made --- which is observed empirically~\citep{zou2023representation} to be the upper-middle layers in similar transformers. Averaging across attention layers avoids over-fitting the extracted direction to any single layer's basis while still excluding the early embedding and SSM layers where the residual is dominated by surface features.

\subsection{Staged training schedule}

Hormone routing is intended to act as a \emph{distillation of the base model's own emerging behavioural axes back into its training signal}. We therefore train in two stages.

\textbf{Stage 1 (base pretraining).} The model is trained from random initialisation as a plain hybrid backbone. The hormone routing module is present in every attention layer, but the hormone bank $V$ is initialised to zero and frozen, so the module computes a zero shift regardless of router output. This stage is identical to training the unmodified hybrid baseline.

\textbf{Hormone extraction.} At a designated milestone --- in our experiments after $1$B training tokens --- we run the contrast-pair extraction procedure of Section~\ref{subsec:extraction} on a frozen copy of the Stage~1 model and write the resulting hormone bank $V$ into the model.

\textbf{Stage 2 (hormone-aware continuation).} Training resumes from the Stage~1 checkpoint with $V$ now set to the extracted directions. The routing weights, per-hormone magnitudes, and output gates are still randomly initialised (router) or zero (gate), so the hormone shift is again zero at the start of Stage~2 and grows only as the gates open. Because $V$ is frozen, the geometric meaning of each hormone direction is preserved through Stage~2; only the routing and gating behaviours evolve.

A natural extension --- which we report in Section~\ref{sec:analysis} but do not adopt as our default --- is to re-extract the hormone bank multiple times during Stage~2, allowing the direction vectors themselves to track the model's evolving representations. Our experiments suggest this is helpful but not necessary for the routing layer to provide a measurable benefit.
